\section{AC (Small-Signal) Analysis}
\label{sec:ac-analysis}

\subsection{Motivation and formulation}

The solver described in Section~\ref{sec:circuit-solver} answers a transient question: given a
circuit's exact component values and a starting state, what does every voltage and current do
over time? A second, complementary question -- how does a circuit respond to a sinusoidal
excitation as a function of frequency alone, independent of any particular starting transient --
is answered instead by AC (small-signal) analysis, the technique underlying every Bode plot in
an electronics textbook. Rather than stepping forward in time, AC analysis assumes every
node voltage and branch current is already varying sinusoidally at a single angular frequency
$\omega$ and solves directly for each phasor's complex amplitude, using exactly the same
node/branch structure as the transient MNA formulation but with every real-valued conductance,
capacitance-derived companion model, and source replaced by its frequency-domain counterpart.

CircuitCraft implements this as an entirely separate solver, \texttt{AcCircuit}, alongside
\texttt{Circuit} rather than as a mode of it: the two share the same node-numbering topology
(Section~\ref{sec:architecture}'s union--find partition is computed once and reused by both,
guaranteeing they agree on which physical position corresponds to which node index) but operate
on disjoint representations -- real-valued matrices advanced tick by tick for the transient
solver, complex-valued matrices solved fresh at each swept frequency for the AC solver -- since
attempting to unify them behind one generic numeric representation would have complicated both
without a compensating benefit.

Following the standard small-signal convention, every independent voltage source in the network
\emph{other than} the one component designated as the AC excitation (Section~\ref{sec:ac-source})
is stamped at exactly zero volts -- an ideal short -- for the duration of the sweep, regardless of
that source's own DC value or redstone-activated state. This isolates the frequency response to
the one signal path actually being characterized, the same simplification a bench engineer makes
by treating every other supply rail as an AC ground when measuring a single stage's response.

\subsection{The AC case for each element}

Every element already described in Section~\ref{sec:circuit-solver} gains a second stamping
method for the AC solver, given the angular frequency $\omega$ directly rather than a timestep:
\begin{itemize}
\item \textbf{Resistor:} impedance is frequency-independent, $Z = R$, so the AC admittance is
simply the same real conductance $1/R$ already used for the transient stamp.
\item \textbf{Capacitor:} $Z = 1/(j\omega C)$, giving a purely imaginary admittance
$Y = j\omega C$ -- no companion history term is needed at all, since AC analysis solves the
steady-state response directly rather than integrating forward through time.
\item \textbf{Inductor:} $Z = j\omega L$, the dual of the capacitor's case, giving
$Y = 1/(j\omega L)$.
\item \textbf{Diode:} linearized about the same last-solved operating point voltage the
transient companion stamp already tracks (Section~\ref{sec:circuit-solver}), into the classic
small-signal diode resistance $r_d = 1/(\mathrm{d}I/\mathrm{d}V) = V_T/I_D$ -- the same
conductance value the transient stamp computes each tick, simply reused as a real admittance
rather than paired with a Norton current source.
\item \textbf{Memristor:} for the moment, treated as a plain resistor frozen at whatever
resistance $R(q)$ it currently holds, with no frequency dependence of its own modeled at all.
This is a deliberate simplification rather than an oversight (Section~\ref{sec:limitations}) --
a genuine small-signal memristor model would need its own frequency response derived from the
charge-controlled state equation, which is left as future work.
\end{itemize}
Each of these was checked directly against a closed-form reference during development
(Section~\ref{sec:validation}): a resistive divider's response is exactly flat across many
decades of frequency; a series RC low-pass reproduces the textbook $-3$~dB, $-45^\circ$ point at
its cutoff frequency $f_c = 1/(2\pi RC)$ and the correct asymptotic $-20$~dB/decade rolloff well
above it; the dual RC high-pass reproduces the complementary $-3$~dB, $+45^\circ$ point; and an
RL divider reproduces the same $-3$~dB behavior at the frequency where the inductor's impedance
magnitude equals the series resistance.

\subsection{The AC Source}
\label{sec:ac-source}

A new two-terminal component, wired like any other (Section~\ref{sec:minecraft-mechanics}),
supplies the sweep's excitation. Unlike every other adjustable component in the mod, none of its
three parameters -- amplitude, and the swept frequency range's minimum and maximum -- are set by
right-click preset cycling: a frequency \emph{range} is not the kind of quantity a short, fixed
preset list represents well, so all three are instead set directly through the shift-right-click
value-editor screen (a sign-style text-entry interface, clamped both client- and server-side to
sensible bounds for each parameter). In the
ordinary transient/DC simulation, the AC Source contributes nothing but a 0~V source -- electrically
a plain wire, exactly like the Ammeter (Section~\ref{sec:architecture}) -- so a circuit built
around one behaves completely normally when nobody is actively running a sweep; its real behavior
appears only through the dedicated probe described next.

\subsection{The two-pole op-amp AC model}

The ideal op-amp of Section~\ref{sec:circuit-solver} is, by construction, an infinite-bandwidth
nullor -- an exact model at the frequencies the transient solver's own worked experiments operate
at, but one that could never exhibit the gain rolloff and phase shift a Bode plot exists to show.
The AC solver therefore replaces it with a standard two-pole open-loop gain model,
\begin{equation}
A(s) = \frac{A_0}{\left(1 + \dfrac{s}{\omega_{p1}}\right)\left(1 + \dfrac{s}{\omega_{p2}}\right)},
\qquad s = j\omega,
\end{equation}
the textbook first-order approximation of a real operational amplifier's frequency response: a
dominant pole from internal frequency compensation, followed by a second, much higher, non-dominant
pole from a secondary internal stage. Fixed, illustrative values are used -- $A_0 = 10^5$ (100~dB
of open-loop DC gain), $f_{p1} = 20$~Hz, $f_{p2} = 3$~MHz -- chosen to make both the dominant
rolloff and the second pole's effect on phase margin visible within a single sweep spanning a few
decades either side, rather than to match any specific real part number. The op-amp remains
modeled the same structural way as its ideal DC/transient counterpart: a voltage-controlled
voltage source with its own branch-current unknown, infinite input impedance, and zero output
impedance, differing only in that the constraint row's coefficient is the complex, frequency-dependent
$A(s)$ rather than a fixed value of~1. This was verified directly: the open-loop gain at a very
low sweep frequency matches $A_0$ to within numerical precision, and the gain measured exactly at
$f_{p1}$ is $3.01$~dB below that DC value, the defining property of a pole.

\subsection{The AC probe and Bode plot}

A second Bode-plot-dedicated probe item captures the specifically two-step interaction this
analysis requires, distinct from every other probe in the mod: right-clicking an AC Source first
only pins it as the sweep's excitation and prints a hint prompting the next step, since nothing
can be computed yet without also knowing which point to observe; right-clicking a second,
different position afterward -- any wired component, wire, or ground block -- runs the sweep and
displays the result. The source stays pinned after that second click, so a player can probe
several different points against the same excitation without re-pinning it each time; only a
shift-right-click, or clicking a different AC Source block, changes it. The sweep itself
evaluates 60 log-spaced frequencies across the pinned source's configured range, solving a fresh
\texttt{AcCircuit} at each one and reporting the magnitude (in dB) and phase (in degrees) of the
signal position's voltage relative to the source's own amplitude -- the standard definition of a
transfer function -- rendered as two stacked graphs against a shared log-frequency axis, the
same layout convention an oscilloscope's own Bode-plot mode uses.
